% ============================================================================
% Fichier  : partie5/P05_C19_audit_continu_outils.tex
% Titre    : Audit Continu, Analyse de Données et Outils
% Auteur   : MINKA MI NGUIDJOÏ Thierry Emmanuel
% Version  : 1.0 -- octobre 2026
% ============================================================================
% Sources : annexe « outils d'automatisation et d'assistance à l'audit »
% du cours Audit_SI de l'auteur (six critères de choix, familles d'outils,
% arbre de décision), reprise et allégée des noms commerciaux (offres et
% éditeurs à vérifier au moment du choix) ; GUID 5100 de l'INTOSAI ;
% article A0 (hal-05701007), thématique T5 « instrumentation adaptative »
% et statut de l'IA générative ; ARKAN : article en préparation, non cité
% en bibliographie. Loi de Benford : résultat classique, présenté comme
% test de détection, sans valeur de preuve.
% ============================================================================

\chapter{Audit Continu, Analyse de Données et Outils}
\label{chap:audit-continu}
\chaptermark{Audit continu et outils}

\epigraph{%
    « L'outil élargit ce que l'auditeur peut regarder.\\
    Il ne lui dit pas ce qu'il faut voir. »%
}{Adage des analystes de données}

\niveauChapitre{Expert}{%
    Maîtrise du Chapitre~\ref{chap:audit-applications} (tests sur les
    données) et du Chapitre~\ref{chap:preuve} (extractions et preuves)%
}

\begin{boxobjectifs}
\begin{itemize}
    \item Appliquer les techniques d'audit assistées par ordinateur, de
          façon reproductible.
    \item Choisir un outil selon six critères, dont l'adaptation au
          contexte africain.
    \item Distinguer contrôle continu et audit continu, et concevoir un
          dispositif pour LiZbiZ.
    \item Présenter les grilles d'audit adaptatives de la thématique T5 de
          l'auditologie.
\end{itemize}
\end{boxobjectifs}

\bigskip

% ============================================================================
\section*{Question directrice : faut-il attendre la prochaine mission~?}
% ============================================================================

Au Chapitre~\ref{chap:audit-applications}, AUDIX CONSEIL a découvert, sur
les données de 2025, six doublons, un fractionnement et quatorze
expéditions non facturées. Le conseil de LiZbiZ pose une question
légitime~: « Pourquoi l'apprendre une fois par an~? Ne peut-on pas le voir
chaque mois~? »

Ce chapitre présente les techniques qui permettent de travailler sur des
populations entières (section~\ref{sec:c19-taao}), le choix d'un outil
(section~\ref{sec:c19-outils}), le passage de la mission ponctuelle au
dispositif continu (section~\ref{sec:c19-continu}), puis les grilles
adaptatives (section~\ref{sec:c19-adaptatif}).

\separateur

% ============================================================================
\section{Les Techniques d'Audit Assistées par Ordinateur}
\label{sec:c19-taao}
% ============================================================================

\subsection{Une population entière, une procédure reproductible}

Les techniques d'audit assistées par ordinateur (\sigle{TAAO}, en anglais
\textit{CAATs}) appliquent des tests à des fichiers de données~: tri,
filtre, jointure, regroupement, recherche de doublons, de trous de
séquence, d'écarts entre fichiers. Le GUID~5100 de l'INTOSAI (§~6.10) en
énumère les usages (analyse des journaux, rapports d'exception,
comparaison de fichiers, stratification, échantillonnage, doublons,
écarts) et trois avantages~: l'analyse de grands volumes, la répétabilité
des tests, et la documentation horodatée des tests et de leurs
résultats~\cite{intosai2019guid5100}.

\begin{tableaucours}{L{3.6cm}L{4.6cm}L{4.0cm}}{%
    \tableauHeader{Test} &
    \tableauHeader{Ce qu'il cherche} &
    \tableauHeader{Dans ce cours}
}
Doublons & Une même opération enregistrée deux fois & Paiements de
LiZbiZ (Chapitre~\ref{chap:audit-applications}, ACSI-LAB-07) \\
Trous et séquences & Numéros manquants ou répétés & Factures, pièces
comptables \\
Rapprochement de fichiers & Un élément présent dans un fichier, absent de
l'autre & Départs et comptes actifs (D.2.1)~; paie et personnel (ACSI-LAB-06) \\
Zones sous un seuil & Concentration juste sous une limite de validation &
TRANSLOG SERVICES (Chapitre~\ref{chap:fraude}) \\
Stratification & Répartition des montants par tranches & Choix d'un
échantillon ciblé \\
Premier chiffre & Écart à la distribution attendue des premiers chiffres &
Montants saisis à la main \\
\end{tableaucours}
\vspace{-0.8em}
\begin{center}{\small\itshape Tableau~19.1. Tests de base des TAAO}\end{center}

\subsection{Le test du premier chiffre}

Dans beaucoup d'ensembles de montants naturels, qui couvrent plusieurs
ordres de grandeur, le premier chiffre significatif $d$ n'est pas
uniforme~: il suit approximativement la loi de Benford,
\[
    P(d) = \log_{10}\!\left(1 + \frac{1}{d}\right),
    \qquad P(1) \approx 30{,}1~\%, \quad P(9) \approx 4{,}6~\%.
\]
Un excès marqué de montants commençant par 4, comme les 4,9~millions de
TRANSLOG SERVICES, attire l'attention. Le test oriente la recherche~; il
ne prouve rien. Il ne s'applique pas aux montants fixés par une règle
(salaires d'une grille, tarifs), ni aux petits ensembles.

\subsection{Rendre une analyse reproductible}

Une analyse de données n'est une preuve que si un tiers peut la refaire et
retrouver le même résultat (axiome~6). L'auditeur conserve donc la chaîne
complète~: l'extraction et son empreinte (Chapitre~\ref{chap:preuve}), le
script ou la suite d'étapes, la version de l'outil, et le résultat.

\begin{boxmetier}{Auditeur SI}
\textbf{Rendre une analyse de données reproductible}
\begin{enumerate}
    \item Ne jamais travailler sur l'extraction d'origine~: en faire une
          copie, et vérifier son empreinte avant chaque séance.
    \item Écrire les traitements dans un script, ou les consigner étape
          par étape, plutôt que de modifier les données à la main.
    \item Noter la version de l'outil, et les paramètres de chaque test
          (seuils, clés de rapprochement, normalisations).
    \item Conserver le résultat avec son empreinte et la date d'exécution.
    \item Faire refaire un test par un collègue avant de rédiger
          l'observation.
\end{enumerate}
\end{boxmetier}

% ============================================================================
\section{Choisir un Outil}
\label{sec:c19-outils}
% ============================================================================

\subsection{Six critères}

\begin{tableaucours}{L{3.4cm}L{8.8cm}}{%
    \tableauHeader{Critère} &
    \tableauHeader{Questions}
}
1. Périmètre & Analyse de données, gestion de la mission et papiers de
travail, gouvernance et risques, contrôle continu, extraction de formats
difficiles~: lesquels faut-il couvrir~? \\
2. Référentiels & L'outil intègre-t-il les cadres utilisés (COSO, COBIT,
ISO/IEC~27001), ou faut-il les paramétrer~? \\
3. Apprentissage & L'équipe sera-t-elle autonome rapidement~? Outil
individuel ou plateforme institutionnelle~? \\
4. Coût & Licences, formation, maintenance~; existe-t-il une offre
académique ou publique~? \\
5. Capacités TAAO & Formats importés, tests disponibles, et surtout
journal des traitements pour la reproductibilité \\
6. Contexte africain & Fonctionnement hors ligne, lecture des exports des
progiciels locaux (par exemple les systèmes de paie publics), coût par
poste soutenable pour une institution publique \\
\end{tableaucours}
\vspace{-0.8em}
\begin{center}{\small\itshape Tableau~19.2. Six critères de choix d'un
outil d'audit}\end{center}

\subsection{Trois familles}

\begin{itemize}
    \item \textbf{Les logiciels d'analyse de données pour auditeurs}, comme
          ACL ou IDEA~: interface sans programmation, journal des
          traitements, tests intégrés~; coût par poste élevé.
    \item \textbf{Les outils libres}~: SQL, tableur, Python avec ses
          bibliothèques d'analyse~; coût nul, grande souplesse, mais
          compétences de programmation et discipline de documentation.
    \item \textbf{Les plateformes de gouvernance, risques et conformité}~:
          cartographie des risques, auto-évaluations, suivi des
          recommandations~; elles servent à piloter, moins à analyser.
\end{itemize}
Les éditeurs et les offres évoluent vite~: les noms cités ici sont des
exemples, à vérifier au moment du choix.

\begin{boxinfo}
\textbf{Un choix raisonné pour une institution publique africaine.} Pour
une équipe sans budget de licences, la combinaison d'un tableur, de SQL et
d'un langage de script couvre les tests du Tableau~19.1, fonctionne hors
ligne et lit les exports en texte. Le ACSI-LAB-07 l'a montré sur
4~812~paiements avec un simple tableur. L'investissement porte alors sur la
formation et sur la discipline de reproductibilité, non sur les licences.
\end{boxinfo}

% ============================================================================
\section{Du Contrôle Continu à l'Audit Continu}
\label{sec:c19-continu}
% ============================================================================

\subsection{Deux dispositifs, deux responsables}

\begin{tableaucours}{L{2.8cm}L{4.6cm}L{4.8cm}}{%
    \tableauHeader{} &
    \tableauHeader{Contrôle continu} &
    \tableauHeader{Audit continu}
}
Qui & Le management, le contrôle interne (première et deuxième lignes) &
L'audit interne (troisième ligne) \\
Objet & Détecter et corriger les anomalies des opérations & Donner une
assurance sur le fonctionnement des contrôles, en continu \\
Exemple & Chaque mois, la comptabilité examine les doublons potentiels &
Chaque mois, l'audit vérifie que l'examen des doublons a eu lieu, et le
refait sur un échantillon \\
\end{tableaucours}
\vspace{-0.8em}
\begin{center}{\small\itshape Tableau~19.3. Contrôle continu et audit
continu}\end{center}

L'audit continu ne remplace pas le contrôle continu~: si l'audit est seul
à surveiller les opérations, il devient lui-même un contrôle, et perd son
indépendance (Chapitre~\ref{chap:normes}).

\subsection{Le dispositif de LiZbiZ}

\begin{boxscenario}{Opération LIZBIZ. Le tableau de bord mensuel}
\textbf{Proposition d'AUDIX CONSEIL}, à partir des tests de la mission~:
\begin{itemize}
    \item \emph{Doublons potentiels} (fournisseur, montant, référence
          normalisée)~: contrôle interne, chaque mois~; alerte dès un cas.
    \item \emph{Paiements entre 4,5 et 5~millions}~: trésorerie et contrôle
          interne, chaque mois~; alerte au-delà de trois pour un même
          fournisseur.
    \item \emph{Expéditions livrées sans facture}~: direction commerciale,
          chaque semaine~; alerte dès un cas de plus de sept jours.
    \item \emph{Comptes actifs d'agents partis}~: DSI et ressources humaines,
          chaque semaine~; alerte dès un cas.
\end{itemize}
\textbf{Rôle de l'audit interne}~: chaque trimestre, vérifier que chaque
alerte a reçu une suite tracée, et refaire un test sur un mois tiré au
sort.
\end{boxscenario}

\begin{boxmetier}{Contrôleur interne}
\textbf{Tenir un indicateur de contrôle continu}
\begin{enumerate}
    \item Écrire la règle de l'indicateur~: données, test, seuil d'alerte,
          fréquence.
    \item Exécuter le test à la date prévue, même quand rien n'est attendu.
    \item Traiter chaque alerte~: explication, correction ou escalade, avec
          sa date.
    \item Conserver l'historique des exécutions et des suites données~:
          c'est la preuve que l'auditeur demandera.
\end{enumerate}
\end{boxmetier}

% ============================================================================
\section{Les Grilles d'Audit Adaptatives}
\label{sec:c19-adaptatif}
% ============================================================================

\subsection{La thématique T5 de l'auditologie}

L'article A0 décrit six thématiques de recherche. La cinquième,
l'\emph{instrumentation adaptative}, vise un outil qui produit des grilles
d'audit adaptées au profil de l'organisation, intègre les mesures du
cours ($\mu$, $U$, $H$, divergence entre cycles) et s'interface avec les
référentiels~\cite{minka2026auditologie}. Sa question de recherche est
directe~: comment concevoir un audit qui s'adapte aux réponses qu'il
reçoit~? Un outil, ARKAN, est en cours de formalisation dans un article
en préparation~; le cours n'en présente que le principe.

\begin{boxdef}
\termecle{Grille adaptative}\index{grille adaptative}~: grille dont les
questions suivantes dépendent des réponses déjà obtenues~: l'effort
d'audit se déplace vers les sous-domaines où l'incertitude reste la plus
forte, au lieu de se répartir également.
\end{boxdef}

\subsection{Une illustration}

\begin{boxlecture}{Domaine D, grille adaptée}
Au Chapitre~\ref{chap:decomposition}, les vingt questions du domaine~D se
répartissaient également~: quatre par sous-domaine. Pour le cycle
suivant, AUDIX CONSEIL applique une règle simple, à titre
d'illustration~: deux questions au moins par sous-domaine, et les dix
restantes réparties selon le taux d'incertitude constaté.

\medskip
\begin{tableaucours}{L{3.6cm}cccc}{%
    \tableauHeader{Sous-domaine} &
    \tableauHeader{O/N/K} & \tableauHeader{$U$} & \tableauHeader{$H$} &
    \tableauHeader{Questions}
}
D.1 Politique & 4/0/0 & 0,00 & 0,00 & 2 \\
D.2 Cycle de vie & 1/1/2 & 0,50 & 1,50 & 2 + 4 = 6 \\
D.3 Authentification & 4/0/0 & 0,00 & 0,00 & 2 \\
D.4 Privilèges & 2/0/2 & 0,50 & 1,00 & 2 + 4 = 6 \\
D.5 Journalisation & 2/1/1 & 0,25 & 1,50 & 2 + 2 = 4 \\
Total & & & & 20 \\
\end{tableaucours}

\medskip
\textbf{Analyse} (confiance moyenne). À effort constant, la grille interroge
trois fois plus D.2 et D.4 que D.1 et D.3, où rien ne restait incertain.
Le minimum de deux questions par sous-domaine évite qu'un domaine stable
ne cesse d'être observé~: c'est le risque de régression du
Chapitre~\ref{chap:cas-integrateurs}. La règle elle-même est une
hypothèse à tester~: l'auditologie demande de vérifier, sur plusieurs
cycles, qu'elle réduit l'incertitude plus vite qu'une grille fixe.

\medskip
\hfill\lecture{0{,}65}{0{,}25}{1{,}24}
\end{boxlecture}

\begin{boxalert}
\textbf{Violation typique.} Laisser un outil choisir seul les questions
d'une mission, sans que l'auditeur puisse expliquer pourquoi un
sous-domaine a été délaissé.

\smallskip
\textbf{Règle~:} la règle d'adaptation est écrite dans le plan d'audit,
discutée avec le commanditaire, et chaque sous-domaine garde un minimum
d'observation. Une grille que l'auditeur ne sait pas expliquer n'est pas
opposable.
\end{boxalert}

% ============================================================================

\separateur

\begin{boxresume}
\textbf{Ce qu'il faut retenir de ce chapitre~:}
\begin{itemize}
    \item \textbf{Les TAAO} travaillent sur des populations entières~:
          doublons, séquences, rapprochements, seuils, stratification,
          premier chiffre.
    \item \textbf{Une analyse n'est une preuve que reproductible}~:
          extraction et empreinte, script, version, résultat.
    \item \textbf{Six critères} guident le choix d'un outil~; les outils
          libres couvrent les tests de base, à condition de rigueur.
    \item \textbf{Le contrôle continu} appartient au management~;
          \textbf{l'audit continu} vérifie qu'il fonctionne, sans s'y
          substituer.
    \item \textbf{Une grille adaptative} déplace l'effort vers
          l'incertitude~; sa règle est écrite, expliquée, et garde un
          minimum d'observation partout.
\end{itemize}
\end{boxresume}

\bigskip

\begin{boxexo}[title={\ding{45}\quad Exercice~19.1. Quel test~? \niveaubase}]
Indiquez le test du Tableau~19.1 adapté à chaque question.
\begin{enumerate}[(a)]
    \item Des factures de vente ont-elles été supprimées~?
    \item Des agents partis sont-ils encore payés~?
    \item Un fournisseur reçoit-il de nombreux paiements juste sous le seuil
          de double validation~?
    \item Les montants des notes de frais saisies à la main paraissent-ils
          naturels~?
\end{enumerate}
\end{boxexo}

\begin{boxexo}[title={\ding{45}\quad Exercice~19.2. Choisir un outil pour l'ANPP \niveaumoyen}]
La Chambre des comptes veut équiper trois magistrats pour auditer les
fichiers de paie de l'ANPP, exportés en texte, dans des locaux où la
connexion est intermittente.
\begin{enumerate}[(a)]
    \item Appliquez les six critères et proposez une solution.
    \item Quelles règles de reproductibilité imposez-vous aux magistrats~?
    \item Quel test de l'outil retenu reprend le ACSI-LAB-06~?
\end{enumerate}
\end{boxexo}

\begin{boxexo}[title={\ding{45}\quad Exercice~19.3. Une grille adaptative pour le domaine E \niveauexpert}]
Le domaine~E de LiZbiZ (exploitation) compte quatre sous-domaines de cinq
questions, cotés en novembre~2026~: E.1 5/0/0~; E.2 1/2/2~; E.3 1/1/3~;
E.4 1/1/3.
\begin{enumerate}[(a)]
    \item Calculez $U$ et $H$ de chaque sous-domaine.
    \item Répartissez vingt questions pour le cycle suivant, avec la règle
          de la lecture « Domaine D, grille adaptée ».
    \item Proposez une autre règle, fondée sur $H$, et comparez.
    \item Comment vérifieriez-vous, sur plusieurs cycles, qu'une règle est
          meilleure que l'autre~?
\end{enumerate}
\end{boxexo}

\bigskip

\begin{boxinfo}
\textbf{Transition.} Les outils d'analyse s'enrichissent aujourd'hui de
modèles d'intelligence artificielle, et les organisations en déploient
dans leurs propres processus. Le Chapitre~\ref{chap:audit-ia} traite des
deux faces de la question~: auditer l'IA, et auditer avec elle.
\end{boxinfo}
